% chapters/02/section-06.tex
\section{Pustaka dan Mesin: Dua Cara Membangun Permainan}
\label{sec:pustaka-mesin}

Bab ini ditutup dengan pembahasan tentang perbedaan paling menentukan pada
mata kuliah ini: perbedaan antara membangun permainan di atas sebuah pustaka
dan membangunnya di dalam sebuah mesin. Perbedaan itulah yang membagi buku ini
menjadi Blok~I--II dan Blok~III--IV.

\subsection{Perbedaan Mendasar}

Sebuah \textbf{pustaka} (\istilah{library}) adalah kumpulan fungsi dan jenis
data yang dapat Anda panggil kapan saja sesuai kebutuhan. Alur program berada
di tangan Anda; pustaka menunggu sampai dipanggil. Ketika Anda menulis
\code{pygame.image.load} untuk memuat gambar, Anda sedang memakai pustaka
\cite{pygame-docs}.

Sebuah \textbf{mesin} (\istilah{engine}) bekerja dengan cara yang berlawanan.
Mesin memiliki alur programnya sendiri, dan kode Anda dipanggil pada
saat-saat yang telah ditentukan. Anda tidak menulis siklus utama permainan;
mesin sudah menjalankannya, dan Anda mengisi bagian pembaruan serta
penayangan. Unity bekerja dengan cara ini, sehingga kode yang Anda tulis
diletakkan di dalam kelas yang dipanggil mesin pada tahap tertentu
\cite{unity-leveldesign,unity-learn}.

Cara kerja mesin semacam itu kadang dirumuskan secara ringkas sebagai
``jangan panggil kami, kami yang akan memanggil Anda''. Rumusan tersebut
menangkap inti perbedaannya: kendali atas alur program berpindah dari kode Anda
ke mesin.

\subsection{Pekerjaan yang Diambil Alih Mesin}

Mesin permainan bukan sekadar pustaka yang lebih besar. Ia mengambil alih
sejumlah pekerjaan yang pada Blok~I--II Anda kerjakan sendiri.

\begin{itemize}
  \item \textbf{Siklus utama permainan}, termasuk pemisahan tahap pembaruan
        berkala dan tahap fisika berlangkah tetap.
  \item \textbf{Penayangan}, termasuk penyusunan urutan gambar, penanganan
        kedalaman, dan penyesuaian terhadap ukuran jendela.
  \item \textbf{Perhitungan fisika}, meliputi gravitasi, tumbukan, gesekan, dan
        tanggapan atas tabrakan.
  \item \textbf{Penyuntingan visual}, sehingga level, posisi, dan ukuran objek
        dapat diatur tanpa menulis kode.
  \item \textbf{Pengelolaan aset}, mencakup impor berkas gambar dan bunyi
        beserta pengaturan penggunaannya.
  \item \textbf{Pemutaran animasi dan bunyi}, termasuk peralihan antar-animasi
        dan pencampuran beberapa sumber bunyi.
  \item \textbf{Penyiapan berkas terpasang} untuk berbagai sistem operasi.
\end{itemize}

\subsection{Konsekuensi bagi Cara Belajar dan Bekerja}

Pengambilalihan tersebut menghasilkan pertukaran yang perlu Anda sadari.

\begin{description}
  \item[Yang diperoleh] Kecepatan pengembangan. Pekerjaan yang pada Blok~I
        menuntut puluhan baris kode dapat dilakukan dalam beberapa menit
        melalui penyuntingan visual. Perhatian dapat dipusatkan pada desain
        level dan pengalaman bermain.
  \item[Yang dikorbankan] Kejelasan sebab akibat. Ketika perilaku permainan
        tidak sesuai harapan, penyebabnya dapat berada pada kode Anda, pada
        pengaturan komponen, atau pada cara mesin bekerja. Membedakan
        ketiganya menuntut pemahaman tentang apa yang sebenarnya dikerjakan
        mesin \cite{unity-leveldesign}.
  \item[Kurva belajar] Mesin menuntut pemahaman tentang cara kerjanya sebelum
        dapat dimanfaatkan sepenuhnya, dan banyak hal harus dipelajari lebih
        dahulu sebelum hasil pertama terlihat.
\end{description}

Di sinilah letak alasan buku ini memakai kedua pendekatan. Pekerjaan pada
Blok~I--II membekali Anda dengan pengetahuan tentang apa yang terjadi di balik
layar: bagaimana waktu dikelola, bagaimana gerak dihitung, dan bagaimana
keputusan NPC disusun. Pengetahuan itu tidak menjadi sia-sia ketika Anda
berpindah ke mesin, melainkan berubah menjadi kemampuan mengenali kapan
pengaturan bawaan sudah memadai dan kapan Anda perlu menulis kode sendiri.
Perbandingan pengalaman bekerja dengan kedua pendekatan akan menjadi bahan
refleksi pada Bab~XVII.

\subsection{Padanan Konsep}

Meskipun berbeda cara kerjanya, kedua dunia itu memecahkan persoalan yang
sama. Tabel~\ref{tab:padanan-konsep} memperlihatkan padanan konsepnya. Tabel
ini berguna sebagai jembatan ketika Anda berpindah dari Blok~II ke Blok~III:
konsep yang telah Anda pahami tetap sama, hanya alat dan letak tanggung
jawabnya yang berubah.

\begin{table}[htbp]
\centering
\caption{Padanan konsep antara Python/Pygame dan Unity}
\label{tab:padanan-konsep}
{\setlength{\tabcolsep}{4pt}\footnotesize%
\begin{tabularx}{\textwidth}{@{}L{3.5cm} Y Y@{}}
\toprule
\headrow \textbf{Konsep} & \textbf{Python dan Pygame} & \textbf{Unity} \\
\midrule
Siklus utama & Perulangan yang Anda tulis sendiri (Bab~IV) &
Dijalankan mesin; Anda mengisi tahap pembaruan dan penayangan (Bab~XIII). \\
\addlinespace
Objek permainan & Objek Python atau kelas turunan \code{Sprite} (Bab~VII) &
\istilah{GameObject} beserta komponen yang menempel padanya (Bab~XII). \\
\addlinespace
Pembaruan waktu & Selang waktu diukur sendiri, lalu dipakai sebagai pengali
(Bab~V) & Disediakan mesin melalui \code{Time.deltaTime} (Bab~XIII). \\
\addlinespace
Gerak dan sifat fisik & Dihitung sendiri dengan vektor dan kaidah Euler
(Bab~V--VI) & Komponen \istilah{Rigidbody} yang perhitungannya dilakukan mesin
(Bab~XIII). \\
\addlinespace
Bentuk batas dan tabrakan & \code{Rect} dan \istilah{mask} (Bab~VII) &
\istilah{Collider2D} atau \istilah{Collider3D} (Bab~XIII). \\
\addlinespace
Level dan peta & Data peta di dalam kode atau berkas tersendiri (Bab~IX--X) &
\istilah{Scene} beserta \istilah{Tilemap} dan \istilah{Prefab}
(Bab~XII--XIII). \\
\addlinespace
Aset gambar & Dimuat saat program dijalankan (Bab~VII) &
Diimpor lebih dahulu, dengan pengaturan impor yang dapat disesuaikan
(Bab~XIV). \\
\addlinespace
Bunyi & Pemutar bunyi di dalam kode (Bab~XV) &
Komponen \istilah{AudioSource} beserta \istilah{AudioClip} (Bab~XV). \\
\bottomrule
\end{tabularx}}
\end{table}

\begin{catatabox}
Kolom \istilah{Konsep} pada tabel tersebut seluruhnya telah
Anda kenal dari bab ini. Yang berubah hanyalah perkakas yang menyediakannya.
Karena itu, apabila pada Blok~III Anda merasa menghadapi dunia yang sama
sekali baru, kembalilah ke bab ini: sebagian besar persoalannya sama, hanya
cara menyelesaikannya yang lebih singkat.
\end{catatabox}
